% Complément à la trame existante. Données archivées, sans nouvelles simulations.
\section{Réponses au sujet et résultats}

\begin{frame}{Le projet : une paroi ferroviaire à quantité limitée}
\begin{itemize}
\item Application retenue : réduire le champ acoustique d'un train dans un tunnel réfléchissant, sur $500$--$3000$ Hz.
\item Produit candidat : support rigide préfractal avec inclusions de laine de roche ; géométrie et placement du matériau sont deux choix distincts.
\item Objectif calculé : $J_f=\int_\Omega |u_f|^2\,dx$, approché sur les nœuds intérieurs. Ce n'est ni la puissance dissipée ni une mesure en dB chez un riverain.
\item Stratégie : identifier les résonances, optimiser le placement, binariser, puis vérifier le \textbf{même mur} sur toute la bande.
\end{itemize}
\medskip
\small Le nom de la startup, les rôles et le coût industriel restent à définir par l'équipe. Le cas autoroutier n'a pas été simulé ici.
\end{frame}

\begin{frame}{État de l'art : géométrie, localisation et absorption}
\begin{itemize}
\item Le support de M. Filoche, \emph{Waves next to irregular boundaries}, relie les frontières irrégulières à la localisation des modes, à l'absorption et aux murs antibruit.
\item Il distingue l'énergie stockée $E_n$ et les pertes par cycle $W_n$ par le facteur de qualité $Q_n=2\pi E_n/W_n$ : les pics et leur largeur renseignent sur les résonances.
\item Notre étude exploite cette piste : modifier une géométrie et la répartition des pertes pour réduire le champ forcé. Une fractale plus longue n'est pas automatiquement meilleure.
\item Les résonateurs constituent une autre piste étudiée localement ; les résultats à deux fréquences doivent être contrôlés hors de leurs cibles.
\end{itemize}
\medskip
\small Appui local : support Filoche, pages 8--12. Cette synthèse ne remplace pas une revue bibliographique des produits, performances expérimentales et coûts.
\end{frame}

\begin{frame}{Le problème de Helmholtz et le budget de matériau}
\small Avec la convention temporelle $e^{-i\omega t}$ et $k=\omega/c$ :
\[
 \Delta u+k^2u=s_f\text{ dans }\Omega,\quad
 u=0\text{ sur }\Gamma_D,\quad \partial_nu=0\text{ sur }\Gamma_N,
 \quad \partial_nu+\alpha(f)\chi u=0\text{ sur }\Gamma_R.
\]
\[
 \min_{\chi\in\mathcal A_Q}J_f(\chi),\qquad
 \mathcal A_Q=\{0\leq\chi\leq1,\ \int_{\Gamma_R}\chi\,ds=Q\}.
\]
\begin{itemize}
\item La relaxation continue permet la descente ; le produit final demande $\chi\in\{0,1\}$.
\item Une fraction $v=30\,\%$ impose $Q_h=h\sum\chi_i=0{,}30hS$ pour $L=1$ m.
\item Quand la paroi s'allonge, \textbf{la même fraction n'impose pas la même quantité}. Comparer les formes exige un $Q_h$ commun.
\item Les essais anciens à source imposée au bord utilisent $u=g\neq0$ en haut : leurs énergies ne sont pas directement comparables au scénario train.
\end{itemize}
\end{frame}

\begin{frame}{Deux matériaux étudiés, un choix encore conditionnel}
\small
\begin{center}
\begin{tabular}{lrrr}
\toprule Matériau & Porosité & Résistivité & Tortuosité\\
\midrule Laine de roche & 0,96 & 32000 & 1,02\\
Caoutchouc recyclé & 0,42 & 18000 & 1,95\\\bottomrule
\end{tabular}
\end{center}
Résistivité en $\mathrm{Pa\,s\,m^{-2}}$ ; paramètres du modèle fourni.
\begin{itemize}
\item $\alpha(\omega)$ est complexe et dépend de la fréquence. Les tables locales contiennent 1000 pulsations correspondant à 500--4000 Hz ; $\omega=2\pi f$.
\item Le code de référence identifie un coefficient de Robin en rapprochant les réponses des modèles volumique et surfacique. Ici, les tables sont réutilisées, sans recalibrage indépendant.
\item À 500 Hz, paroi plane et source au bord : énergie finale $0{,}498$ pour le caoutchouc contre $0{,}658$ pour la laine ($N=128$, $Q=0{,}4$).
\item Sur le diagnostic 500--899,4 Hz des quatre formes, les pics finaux sont plus faibles avec la laine. Sa supériorité universelle n'est pas démontrée.
\end{itemize}
\end{frame}

\begin{frame}{Un mur fixé au pic initial : protocole des figures jointes}
\begin{enumerate}
\item Balayer la bande pour la répartition initiale et la couverture intégrale.
\item Choisir le maximum initial échantillonné : $1043{,}0$ Hz pour $j=1$, $706{,}7$ Hz pour $j=2$.
\item Optimiser une seule répartition continue à cette fréquence, puis appliquer les deux binarisations.
\item Recalculer les spectres des trois répartitions \textbf{fixées}, sans réoptimiser ailleurs.
\end{enumerate}
\medskip
\small Figures jointes : laine de roche, source train, 714 fréquences. Les métadonnées complètes de ces six PNG ne sont pas jointes : ne pas déduire le maillage ou la quantité après seuil des seules courbes.
\medskip

La couverture à $100\,\%$ signifie $\chi=1$ partout : \textbf{elle ne signifie pas une absorption acoustique parfaite}.
\end{frame}

\begin{frame}{Niveau 1 : le pic ciblé baisse, un autre devient dominant}
\centering
\includegraphics[width=.94\textwidth,height=.65\textheight,keepaspectratio]{figures_presentation/j1_Figure_2.png}

\small Le maximum passe de $2{,}89\,10^{-4}$ à $3{,}10\,10^{-4}$ (seuil)
ou $3{,}27\,10^{-4}$ (tri), vers $1719{,}2$ Hz : \textbf{environ $+7{,}3\,\%$ et $+13{,}1\,\%$}.
\end{frame}

\begin{frame}{Niveau 2 : le maximum spectral diminue également}
\centering
\includegraphics[width=.94\textwidth,height=.65\textheight,keepaspectratio]{figures_presentation/j2_Figure_2.png}

\small Le maximum passe de $2{,}50\,10^{-4}$ à $8{,}39\,10^{-5}$ (seuil)
ou $1{,}17\,10^{-4}$ (tri) : \textbf{environ $-66{,}4\,\%$ et $-53{,}2\,\%$}.
Le maximum restant est à $1719{,}2$ Hz.
\end{frame}

\begin{frame}{La somme sur la bande baisse pour les deux niveaux}
\small Valeurs lues sur les histogrammes joints ; unités du modèle.
\begin{center}
\begin{tabular}{lrrrr}
\toprule & Initial & Seuil & Tri & Couverture 100\,\%\\
\midrule $j=1$ & $1{,}280\,10^{-3}$ & $1{,}011\,10^{-3}$ & $1{,}040\,10^{-3}$ & $1{,}615\,10^{-3}$\\
Variation & référence & $-21{,}0\,\%$ & $-18{,}8\,\%$ & $+26{,}2\,\%$\\
\addlinespace
$j=2$ & $1{,}194\,10^{-3}$ & $6{,}899\,10^{-4}$ & $8{,}108\,10^{-4}$ & $1{,}470\,10^{-3}$\\
Variation & référence & $-42{,}2\,\%$ & $-32{,}1\,\%$ & $+23{,}1\,\%$\\\bottomrule
\end{tabular}
\end{center}
\begin{itemize}
\item Les pourcentages reprennent les annotations arrondies des figures.
\item Le seuil donne ici la plus faible somme ; sa quantité de matériau doit être vérifiée avant de conclure à une meilleure efficacité à budget égal.
\item Une somme $\sum_f J_f$ dépend de l'échantillonnage. Ce n'est pas une énergie acoustique totale physique ; une intégrale fréquentielle requiert des poids de quadrature.
\end{itemize}
\end{frame}

\begin{frame}{Pourquoi couvrir toute la paroi ne minimise-t-il pas toujours $J$ ?}
\begin{itemize}
\item Le coefficient complexe de Robin modifie à la fois les pertes et la réponse résonante du domaine.
\item Changer $\chi$ change le champ $u$ : davantage de matériau n'implique pas une décroissance point par point de $\int|u|^2$.
\item Dans les deux essais joints, la couverture intégrale a une somme plus élevée que la répartition initiale. Ce constat ne prouve pas qu'elle dissipe moins de puissance.
\item Les maxima peuvent se déplacer entre les points calculés : les étoiles désignent des \textbf{maxima échantillonnés}.
\end{itemize}
\medskip
\textbf{Il faut évaluer séparément le champ, les flux dissipés et les pics.}
L'optimisation continue monotone à une fréquence ne garantit ni le meilleur binaire, ni le meilleur mur sur la bande.
\end{frame}

\begin{frame}{Optimisations successives : traiter le pic qui reste}
\small Archive \code{resultats_optimisation_pics_successifs} :
$j=2$, laine de roche, $N=128$, 714 fréquences, six cycles, 80 itérations maximum par cycle, source train à 300 km/h.
\begin{enumerate}
\item Détecter le maximum du spectre de la répartition courante.
\item Optimiser à cette fréquence, puis binariser.
\item Recalculer tout le spectre et utiliser cette paroi pour le cycle suivant.
\end{enumerate}
\medskip
\textbf{Ce n'est pas une optimisation simultanée multifréquence.}
Un pic déjà traité peut réapparaître et la somme peut remonter.
\medskip

\scriptsize Traçabilité : le JSON archivé indique binarisation 1, seuil $0{,}5$ ;
le script actuel appelle \code{chi_binarize_2}. Les résultats présentés sont ceux de l'archive, pas d'une nouvelle exécution du code actuel. La fraction archivée après binarisation vaut $29{,}941\,\%$.
\end{frame}

\begin{frame}{Les cycles successifs ne s'améliorent pas de façon monotone}
\centering
\includegraphics[width=.96\textwidth,height=.68\textheight,keepaspectratio]{resultats_optimisation_pics_successifs/sommes_energies.png}

\small Le cycle 4 minimise la somme parmi les six murs binaires :
$7{,}0074\,10^{-4}$, soit \textbf{$41{,}33\,\%$ de baisse} par rapport à l'état initial.
\end{frame}

\begin{frame}{Le meilleur cycle dépend du critère choisi}
\small Valeurs recalculées depuis \code{energies.csv} ; maximum sur les 714 points.
\begin{center}
\begin{tabular}{rrrr}
\toprule Cycle & Somme ($10^{-4}$) & Maximum ($10^{-5}$) & Fréquence du max. (Hz)\\
\midrule
0 & 11,9440 & 24,9870 & 706,7\\
1 & 8,3237 & 13,9873 & 1719,2\\
2 & 8,8361 & 6,4328 & 1043,0\\
3 & 7,7248 & 3,4393 & 1719,2\\
4 & 7,0074 & 4,6918 & 1715,7\\
5 & 7,6753 & 9,1957 & 1719,2\\
6 & 7,1119 & 5,0247 & 706,7\\
\bottomrule
\end{tabular}
\end{center}
\textbf{Cycle 3 : meilleur maximum ; cycle 4 : meilleure somme.}
\medskip

L'essai joint $j=2$ avec seuil affiche une somme $6{,}899\,10^{-4}$, légèrement plus basse que celle du cycle 4. Les cycles ne démontrent donc pas une supériorité systématique ; les paramètres et quantités doivent d'abord être alignés.
\end{frame}

\begin{frame}{Le mur retenu selon la somme conserve des pics résiduels}
\centering
\includegraphics[width=.96\textwidth,height=.68\textheight,keepaspectratio]{resultats_optimisation_pics_successifs/comparaison_initial_meilleur_spectre.png}

\small Même répartition du cycle 4 à toutes les fréquences. Son maximum vaut
$4{,}692\,10^{-5}$ vers $1715{,}7$ Hz : baisse d'environ $81{,}2\,\%$ par rapport au maximum initial de cette archive.
\end{frame}

\begin{frame}{Que disent les essais antérieurs sur les formes ?}
\small Diagnostic distinct : source plane au bord, $Q=0{,}4$, $N=128$,
optimisation simultanée à $500$ et $612{,}112$ Hz, puis mur fixe sur 500--899,4 Hz.
\begin{center}
\begin{tabular}{lrr}
\toprule Forme & Pic final laine & Pic final caoutchouc\\
\midrule Plane & 2,886 & 21,124\\
Fractale niveau 1 & 3,613 & 33,276\\
Plane + résonateur & 2,543 & 22,131\\
Fractale + résonateur & 3,921 & 18,062\\\bottomrule
\end{tabular}
\end{center}
\begin{itemize}
\item La forme plane avec résonateur en laine donne le plus faible pic de ce diagnostic, $11{,}9\,\%$ sous la référence plane en laine.
\item Le mélange fractale + résonateur est meilleur sur les deux fréquences cibles, mais pas sur le pire pic de la sous-bande.
\item Ne pas comparer ces valeurs absolues avec celles de la source train ; les excitations diffèrent. Le balayage large bande reste nécessaire.
\end{itemize}
\end{frame}

\begin{frame}{Ce que l'on peut proposer comme produit}
\begin{itemize}
\item \textbf{Piste ferroviaire :} paroi préfractale de niveau 2 avec inclusions de laine de roche. Les résultats joints justifient de poursuivre ce candidat, sans établir un optimum de forme.
\item Choisir le mur selon le cahier des charges : somme sur la bande ou pire résonance. Pour l'archive des cycles, ces choix donnent respectivement les cycles 4 et 3.
\item Conserver une paroi plane et une variante à résonateur comme références à quantité totale identique.
\item Avant fabrication : confirmer le maillage, le budget, la robustesse aux sources et les performances dans une configuration ouverte.
\end{itemize}
\medskip
\small Aucun coût, gain mesuré en dB, vieillissement, comportement mécanique ou résistance aux intempéries n'est évalué. Le choix est un \textbf{candidat numérique}, pas une validation industrielle.
\end{frame}

\begin{frame}{Les calculs complémentaires les plus utiles}
\small
\begin{enumerate}
\item \textbf{Optimisation simultanée des pics} : même $\chi$ à 706,7, 1043,0 et 1719,2 Hz ; ajouter les nouveaux pics détectés et comparer somme/max au meilleur cycle.
\item \textbf{Convergence} : fixer une paroi physique, comparer $N=128,256$ et raffiner autour de 1715--1725 Hz. Valider les pics avant de classer les murs.
\item \textbf{Quantité optimale} : tracer le compromis énergie/quantité pour plusieurs $Q$, avec tri à budget fixé et mêmes sources/fréquences pour toutes les formes.
\item \textbf{Sources indépendantes et frontière ouverte} : séparer roulement/aérodynamique ; remplacer le bord supérieur par une condition rayonnante et comparer les flux.
\end{enumerate}
\medskip
Les scripts existants couvrent une partie de ces besoins ; le fichier
\code{CALCULS_COMPLEMENTAIRES_EI.md} précise les relances et les développements nécessaires.
\end{frame}

\begin{frame}{Couverture du sujet : bilan honnête}
\small
\begin{tabular}{p{.25\textwidth}p{.65\textwidth}}
\toprule Demande & Réponse disponible et limite\\\midrule
Jour 1 & Modèle, sources, matériau, contrainte et gradient présentés ; théorie en annexe. État de l'art local partiel ; organisation à compléter.\\[1mm]
Jour 2 & Pics et optimisation ciblée pour $j=1,2$ ; comparaison à $\chi=1$. Étude plane intégralement couverte et comparaison exhaustive des formes encore à consolider pour la source train.\\[1mm]
Jour 3 & Mur fixe sur la bande et cycles analysés. Multifréquence simultané réalisé dans les anciens essais à deux cibles, pas pour les pics du train. Quantité optimale non établie.\\[1mm]
Jour 4 & Candidat proposé, limites et validations identifiées.\\[1mm]
Jour 5 & Support complété ; sélectionner les diapositives pour 20 minutes, répartir les interventions et finaliser le rapport d'équipe.\\\bottomrule
\end{tabular}
\end{frame}

\appendix
\section{Annexes : théorie et vérifications}

\begin{frame}{Formulation faible et propriétés de la solution}
\small Posons $V=\{v\in H^1(\Omega):v|_{\Gamma_D}=0\}$. Pour la source volumique :
\[
 a_\chi(u,v)=\int_\Omega \nabla u\cdot\nabla\bar v-k^2u\bar v\,dx
 +\int_{\Gamma_R}\alpha\chi u\bar v\,ds
 =-\int_\Omega s_f\bar v\,dx.
\]
\begin{itemize}
\item Sur les préfractales polygonales finies, les traces usuelles sont définies. La forme est continue mais non coercive en général à cause de $-k^2$.
\item Avec absorption passive ($\operatorname{Im}\alpha<0$) sur une portion ouverte où $\chi>0$, l'identité imaginaire du problème homogène impose $u=0$ sur cette portion ; Robin impose aussi $\partial_nu=0$.
\item Sous les hypothèses usuelles de régularité et de continuation unique, cela donne l'unicité ; l'alternative de Fredholm donne l'existence et la dépendance continue aux données.
\item Sans absorption, des fréquences propres peuvent compromettre l'unicité. Près d'une résonance, la réponse et le conditionnement peuvent être élevés.
\end{itemize}
\scriptsize Pour une source de Dirichlet non homogène, effectuer un relèvement $u=u_g+w$. Aucun théorème pour une fractale limite n'est nécessaire aux géométries finies simulées ici.
\end{frame}

\begin{frame}{Dérivée de l'énergie et état adjoint}
\small À géométrie et source fixées, pour une variation réelle $\eta$ de $\chi$ :
\[
 a_\chi(u',v)=-\int_{\Gamma_R}\alpha\eta u\bar v\,ds,
 \qquad J'(\chi)\eta=2\operatorname{Re}\int_\Omega u'\bar u\,dx.
\]
Définir l'adjoint $p\in V$ par
\[
 a_\chi(v,p)=\int_\Omega v\bar u\,dx\quad\forall v\in V.
\]
Alors
\[
 \boxed{J'(\chi)\eta=-2\operatorname{Re}\int_{\Gamma_R}\alpha u\bar p\,\eta\,ds.}
\]
Dans le code discret, $A^Hp=h^2Mu$, avec $M$ masque des nœuds intérieurs :
\[
 dJ=2\operatorname{Re}\big[p^H(db-dA\,u)\big].
\]
\scriptsize La dérivée doit inclure l'élimination des conditions de Robin dans $A$ et $b$. Une différence finie directionnelle permet de vérifier l'implémentation. Ici, $p$ n'inclut ni le facteur 2 ni un signe moins.
\end{frame}

\begin{frame}{Frontière ouverte et réfléchissante : bilan énergétique}
\small Pour $\Delta u+k^2u=s$, multiplier par $\bar u$, intégrer et utiliser Green :
\[
 \int_\Omega |\nabla u|^2-k^2|u|^2\,dx+
 \int_{\Gamma_R}\alpha\chi|u|^2\,ds
 -ik\int_{\Gamma_o}|u|^2\,ds=-\int_\Omega s\bar u\,dx.
\]
Sur $\Gamma_o$, on a utilisé $\partial_nu-iku=0$ ; les autres bords sont homogènes.
\[
 \underbrace{-\int_{\Gamma_R}\operatorname{Im}(\alpha)\chi|u|^2\,ds}_{\text{pertes sur la paroi}}
 +\underbrace{k\int_{\Gamma_o}|u|^2\,ds}_{\text{flux sortant}}
 =\operatorname{Im}\int_\Omega s\bar u\,dx.
\]
\begin{itemize}
\item Bord rigide : $\partial_nu=0$, donc pas de flux sortant par ce bord.
\item L'ouverture ajoute une évacuation d'énergie, mais le champ et le travail de la source changent : cette identité ne prouve pas $J_{\text{ouvert}}\leq J_{\text{fermé}}$ à toute fréquence.
\end{itemize}
\scriptsize Flux proportionnels à une puissance après normalisation physique. Sommerfeld locale est une approximation de bord transparent, surtout pour l'incidence normale ; ce n'est pas une représentation exacte de tout l'espace extérieur.
\end{frame}

\begin{frame}{Dériver par rapport à la forme : ce qui reste à faire}
\small Le gradient précédent porte sur la \textbf{répartition du matériau à forme fixée}.
Les comparaisons plane/fractale/résonateur ne sont pas une descente de forme.
\medskip

Pour une déformation $T_t=I+t\theta$ du domaine et un champ différentiable :
\[
 \frac{d}{dt}J(\Omega_t)\big|_{0}
 =2\operatorname{Re}\int_\Omega u'\bar u\,dx
  +\int_{\partial\Omega}|u|^2\theta\cdot n\,ds.
\]
\begin{itemize}
\item $u'$ doit satisfaire l'équation et les conditions aux bords différentiées ; celles-ci diffèrent entre Robin, Neumann et Sommerfeld.
\item Pour une vraie optimisation de forme, il faut fixer le transport de $\chi$, les bords mobiles et la conservation de $Q$, puis dériver le Lagrangien et vérifier sa dérivée.
\item La formule de transport seule n'est pas un gradient de forme exploitable. Cette partie de la demande théorique reste \textbf{partielle}.
\end{itemize}
\end{frame}

\begin{frame}{Optimiser plusieurs fréquences simultanément}
\[
 F(\chi)=\sum_{\ell=1}^m w_\ell J_{f_\ell}(\chi),\qquad
 \nabla F=\sum_{\ell=1}^m w_\ell\nabla J_{f_\ell},\quad w_\ell\geq0.
\]
\begin{enumerate}
\item Résoudre un état direct et un adjoint pour chaque fréquence.
\item Sommer les gradients, descendre puis projeter sur $\mathcal A_Q$.
\item Accepter selon \textbf{l'objectif total}, puis binariser et rebalayer la bande.
\end{enumerate}
\small Pour limiter le pire pic : minimiser $\max_f J_f$, ou une approximation lisse,
plutôt que seulement la moyenne. Enrichir les cibles lorsque de nouveaux pics apparaissent.
\medskip

Les deux sources train sont actuellement cohérentes. Pour des sources indépendantes,
utiliser $J=J_{\rm roul}+J_{\rm aero}$ avec deux résolutions séparées, sans terme d'interférence.
\end{frame}

\begin{frame}{Maillage, convergence et reproductibilité}
\small Le sujet impose au moins dix points par longueur d'onde :
\[
 \lambda/h=\frac{c}{fh}\geq10,\qquad h=1/N.
\]
À 3000 Hz et $c=340$ m/s : $N\geq89$ ; $N=128$ donne $14{,}5$ points par longueur d'onde.
\begin{itemize}
\item Cette condition est nécessaire, mais ne valide ni les détails géométriques ni les pics étroits. Le pas spectral des tables est environ $3{,}5035$ Hz.
\item Les essais anciens $N=64/128$ à 500 Hz présentent encore des écarts d'énergie de 2 à 19\,\%. Le contrôle hybride $N=256$ n'étudie que les états initiaux à deux fréquences.
\item Pour chaque relance : conserver $\chi$, fréquence, $\alpha$, $N$, $Q$, graine éventuelle, historique, méthode binaire et version du code.
\item Les 80 pas acceptés de chaque cycle archivé ne prouvent pas une convergence vers un optimum ; mesurer aussi le gradient projeté.
\end{itemize}
\end{frame}

\begin{frame}{Figures jointes : spectre complet du niveau 1}
\centering\includegraphics[width=.98\textwidth,height=.78\textheight,keepaspectratio]{figures_presentation/j1_Figure_1.png}
\end{frame}
\begin{frame}{Figures jointes : spectre complet du niveau 2}
\centering\includegraphics[width=.98\textwidth,height=.78\textheight,keepaspectratio]{figures_presentation/j2_Figure_1.png}
\end{frame}
\begin{frame}{Figures jointes : somme au niveau 1}
\centering\includegraphics[width=.98\textwidth,height=.78\textheight,keepaspectratio]{figures_presentation/j1_Figure_3.png}
\end{frame}
\begin{frame}{Figures jointes : somme au niveau 2}
\centering\includegraphics[width=.98\textwidth,height=.78\textheight,keepaspectratio]{figures_presentation/j2_Figure_3.png}
\end{frame}

\begin{frame}{Sources et portée des conclusions}
\scriptsize
\begin{itemize}
\item Cahier des charges : \code{Sujet-EI-EXT-2026.pdf}, deux pages.
\item Figures jointes : \code{Downloads/resultats/j1} et \code{j2}, figures 1 à 3. Les valeurs de synthèse proviennent des annotations arrondies.
\item Cycles : \code{resultats_optimisation_pics_successifs/energies.csv}, \code{bilan.json}, figures du même dossier. Les tableaux sont recalculés à partir du CSV.
\item Essais antérieurs : \code{RESULTATS_NIVEAUX.md}, \code{RESULTATS_HYBRIDE.md}, \code{resultats_pics_hybrides/synthese_pics.csv}. Conditions de source distinctes.
\item Matériaux : \code{data_alpha}, \code{references_mauri/compute_alpha.py} et \code{COMPARAISON_MATERIAUX.md}. Tables reprises, pas d'identification indépendante.
\item Appuis théoriques locaux : \code{Dirichlet_Nonhomogene.pdf}, \code{GAllaire-OptimParam.pdf}, \code{Mur-antibruit-Filoche.pdf}. État de l'art à compléter avant remise ; aucune comparaison industrielle exhaustive n'est revendiquée.
\item Implémentation : \code{demo_control_polycopie2024.py}, \code{processing.py}, \code{train_source.py} et les deux tests de spectre/pics.
\end{itemize}
\end{frame}
